\section{Introduction}
Vector search increasingly appears inside structured database queries.
A catalog search combines semantic similarity with product categories and
availability; document retrieval restricts embeddings by topic or access
labels. The index must find nearby vectors among those satisfying the whole
predicate. We study conjunctions of labels: a vector is eligible when its
label set contains every query label. Selectivity denotes the fraction of
eligible vectors.

Graph indexes must reconcile two organizations of this data: proximity in
vector space and eligibility in label space. Filter-aware expansion, as in
ACORN~\cite{patel2024acorn}, improves navigation through a global graph.
Label-aware systems instead expose metadata structure. UNG~\cite{cai2024ung}
partitions vectors by exact label set, connects groups through minimal
strict-superset relations, and realizes those relations as vector edges.
Curator~\cite{jin2026curator} shares label-specialized virtual indexes through
a clustering tree. Our starting point is UNG's containment organization,
which gives every cross-group transition a clear eligibility argument.

That organization couples two sources of work. Search first identifies the
minimal observed label sets containing the query. It then traverses a vector
graph whose group transitions respect containment. Entry minimality may
require comparisons across unrelated label branches, while exact-group
boundaries constrain which geometric transitions are available. On Amazon,
602,453 vectors occupy 482,387 exact-label groups. A query can consequently
spend substantial work both finding entries and navigating fragmented groups.
The two costs need to be addressed together: reducing one can increase the
other.

\paragraph{Design entries and connectivity together.}
A canonical label trie connects each observed label set to its nearest
terminal descendants. Its entry set consists of the first eligible terminal
on every eligible prefix branch. This \emph{prefix frontier} covers all
eligible groups in the resulting forest. It can be found without comparing
eligible terminals across different branches for set minimality. The frontier
may contain more groups than UNG's minimal-superset entries, but those extra
entries are what make the simpler connectivity usable. Figure~\ref{fig:method-overview}
shows this dependency on one catalog query.

\paragraph{Authorize a region, then navigate geometrically.}
The same trie supplies a way to relax exact-group boundaries. Every vector
below a prefix $R_b$ contains $R_b$. If the query labels $Q$ satisfy
$Q\subseteq R_b$, a local proximity graph may connect these vectors across
groups in either label direction. The prefix thus certifies a whole search
region. We materialize such regions as \emph{blocks}, whose direct members
exclude separately emitted descendant blocks. Queries that cannot authorize
a block still need the finer group representation. This yields an exact-group
base and independently partitioned block overlays, with one overlay as the
simplest instance.

\system{} combines these mechanisms through level-tagged search states. Each
state expands the adjacency of one physical level; entry selection and block
authorization determine which states enter the bounded search. This makes
predicate safety independent of the chosen LNG or prefix connectivity at each
level. We also batch irregular graph-construction tasks on the GPU and give a
structural rule that derives overlay thresholds and depth from dataset size
and existing degree budgets.

We evaluate the tradeoff by separating matched entry/topology systems from
controlled upper-level topology changes. The complete Amazon factorial
contains 14 configurations over nine selectivities. At measured Recall@10
$\geq0.90$, the prefix-base system is 2.77--18.43$\times$ faster than the
optimized LNG base on three workloads near 0.5--5\%, whereas LNG wins near
10\% and 30\%. Adding blocks changes this picture again: the best systems
at 60\% and 80\% combine a Trie base with an LNG overlay. The detailed
profiles explain how discovery savings can be offset by seed and graph work.


The contributions are:
\begin{itemize}
\item \textbf{A joint entry and topology design.} A terminal-prefix forest
  and bitmap-backed frontier provide label-plane coverage with an explicit
  tradeoff between minimality work and entry count.
\item \textbf{Predicate-certified block navigation.} Independent prefix
  partitions let queries use coarser vector graphs when a root-prefix
  certificate holds. We specify the seeding and level-local search algorithm
  and establish its conditional safety properties.
\item \textbf{An analysis of the entry--navigation tradeoff.} A complete
  topology factorial and query profiles identify where discovery savings
  are displaced into seed and graph work, and where blocks recover throughput.
  Supporting studies examine structural configuration and GPU construction.
\end{itemize}
